% =====================================================================
% ODD protocol LaTeX template — version 2010 (first update)
%
% Structure and questions below are taken from the OFFICIAL template
% supplement of:
%   Grimm V., Berger U., DeAngelis D.L., Polhill J.G., Giske J.,
%   Railsback S.F. (2010) "The ODD protocol: A review and first update."
%   Ecological Modelling 221(23): 2760-2768.
%   DOI: 10.1016/j.ecolmodel.2010.08.019
%   Green OA PDF of the paper + template supplement:
%   http://www.agsm.edu.au/bobm/teaching/SimSS/ODD_protocol.pdf
%
% How to use this template:
%   - Replace every "\textit{[...]}"" placeholder with your answer.
%   - Delete every shaded "Question" block before submission
%     (official rule: "Questions and explanations should, of course,
%     be deleted then." — ODD template supplement, 2010).
%   - The % comments carry the official explanations/rules; they do
%     not appear in the PDF. Keep them while drafting.
%   - The 7 element titles and their numbering (1-7) must be used
%     exactly as given.
% =====================================================================
\documentclass[11pt,a4paper]{article}

\usepackage[utf8]{inputenc}
\usepackage[T1]{fontenc}
\usepackage[english]{babel}
\usepackage{geometry}
\geometry{margin=2.5cm}
\usepackage{xcolor}
\usepackage{mdframed}

% Style for the official ODD questions (delete blocks before submission)
\definecolor{oddquestioncolor}{gray}{0.35}
\newmdenv[linewidth=0.5pt,linecolor=gray!50,backgroundcolor=gray!8,
  innertopmargin=6pt,innerbottommargin=6pt,leftmargin=8pt,
  rightmargin=8pt,splittopskip=6pt]{oddquestion}
\newenvironment{question}
  {\begin{oddquestion}\small\itshape\color{oddquestioncolor}%
   \textbf{Question (delete this box before submission):}\par\smallskip}
  {\end{oddquestion}}

\title{ODD model description\thanks{Template following the ODD protocol
(Overview, Design concepts, Details) of Grimm et al. (2006, 2010).}}
\author{Author names}
\date{}

\begin{document}
\maketitle

% Opening sentence required by the official template (keep it):
The model description follows the ODD (Overview, Design concepts,
Details) protocol for describing individual- and agent-based models
(Grimm et al.\ 2006, 2010).

% =====================================================================
% CATEGORY: Overview  (comment only — categories are not part of the
% ODD document itself; only the seven numbered elements are)
% =====================================================================

\section{Purpose}
\begin{question}
What is the purpose of the model?
\end{question}
\textit{[State concisely the overall objective(s) for which the model was
developed. Describe only what the model is to be used for, not how it
works. Make this paragraph independent of any purpose statement in the
article's introduction: ODD must be complete and understandable by
itself. If one purpose is to expand from basic principles to richer
representation of real-world scenarios, state this explicitly.]}

% =====================================================================

\section{Entities, state variables, and scales}
\begin{question}
What kinds of entities are in the model? By what state variables, or
attributes, are these entities characterized? What are the temporal and
spatial resolutions and extents of the model?
\end{question}
\textit{[Describe: (a) entities — typically agents/individuals (with
their types), spatial units such as grid cells, the overall environment,
and collectives; (b) state variables/attributes for each entity type,
with units where applicable. State variables must be low-level and
elementary: they cannot be calculated from other state variables
(distance to a site, e.g., is NOT a state variable if positions are).
A useful test: if you had to stop the model and re-start it later in
exactly the same state, what information would you have to save?
(c) temporal and spatial resolution and extent, specifying what the
model's units represent in reality, e.g.: "One time step represents one
year and simulations were run for 100 years. One grid cell represents
1 ha and the model landscape comprised 1,000 x 1,000 ha; i.e., 10,000
square kilometers."]}
% Official explanation (2010 supplement): entities are distinct objects
% or actors that behave as a unit and interact; agents may include
% identity number, age, sex, location, size, weight, energy reserves,
% behavioral strategy, etc. State variables may change over time (weight)
% or remain constant (sex, species parameters).

% =====================================================================

\section{Process overview and scheduling}
\begin{question}
Who (i.e., what entity) does what, and in what order? When are state
variables updated? How is time modeled, as discrete steps or as a
continuum over which both continuous processes and discrete events can
occur?
\end{question}
\textit{[Describe the schedule in pseudo-code, in every detail, so that
the model can be re-implemented from this code; ideally the pseudo-code
corresponds fully to the actual code. List only the self-explanatory
names of processes here ('update habitat', 'move', 'grow', ...); these
names become the titles of the submodels described in Element 7.
Specify, for each process: which entity performs it, the order among
processes, and the order among the entities performing one process
(random, fixed, size-sorted, ... — ordering can strongly affect results).
State whether updating is asynchronous (new value assigned immediately)
or synchronous (values stored and updated at once for all agents), and
how time is modeled if not already clear from Element 2.]}

% =====================================================================
% CATEGORY: Design concepts
% =====================================================================

\section{Design concepts}
\begin{question}
There are eleven design concepts, discussed by Railsback (2001) and
Grimm \& Railsback (2005, Ch.~5), summarized via the questions of each
concept below.
\end{question}
% Official rule (2010 supplement): unused concepts can be omitted from
% the ODD description. If an important concept underlying the model's
% design is NOT part of ODD, give it a short name, clearly announce it
% as a design concept not included in the ODD protocol, and present it
% at the END of this element.
% The "Design concepts" element does not describe the model per se (it
% is not needed to replicate it) but is a checklist ensuring that
% important design decisions are made consciously and made known to
% readers.

\subsection{Basic principles}
\begin{question}
Which general concepts, theories, hypotheses, or modeling approaches
are underlying the model's design? Explain the relationship between
these basic principles, the complexity expanded in this model, and the
purpose of the study. How were they taken into account? Are they used
at the level of submodels (e.g., decisions on land use, or foraging
theory), or is their scope the system level (e.g., intermediate
disturbance hypotheses)? Will the model provide insights about the
basic principles themselves, i.e., their scope, their usefulness in
real-world scenarios, validation, or modification? Does the model use
new, or previously developed, theory for agent traits from which system
dynamics emerge?
\end{question}
\textit{[Answer here.]}

\subsection{Emergence}
\begin{question}
What key results or outputs of the model are modeled as emerging from
the adaptive traits, or behaviors, of individuals? In other words, what
model results are expected to vary in complex and perhaps unpredictable
ways when particular characteristics of individuals or their environment
change? Are there other results that are more tightly imposed by model
rules and hence less dependent on what individuals do, and hence
'built in' rather than emergent results?
\end{question}
\textit{[Answer here.]}

\subsection{Adaptation}
\begin{question}
What adaptive traits do the individuals have? What rules do they have
for making decisions or changing behavior in response to changes in
themselves or their environment? Do these traits explicitly seek to
increase some measure of individual success regarding its objectives
(e.g., ``move to the cell providing fastest growth rate'', where growth
is assumed to be an indicator of success)? Or do they instead simply
cause individuals to reproduce observed behaviors (e.g., ``go uphill
70\% of the time'') that are implicitly assumed to indirectly convey
success or fitness?
\end{question}
\textit{[Answer here.]}

\subsection{Objectives}
\begin{question}
If adaptive traits explicitly act to increase some measure of the
individual's success at meeting some objective, what exactly is that
objective and how is it measured? When individuals make decisions by
ranking alternatives, what criteria do they use? Synonyms for
'objectives' are 'fitness' for organisms assumed to have adaptive traits
evolved to provide reproductive success, 'utility' for economic reward
in social models, or simply 'success criteria'. (The objective of agents
that are members of a team, social insects, organs of an organism, or
cells in a tissue may refer to the team, colony or organism, not to
themselves.)
\end{question}
\textit{[Answer here.]}

\subsection{Learning}
\begin{question}
Many individuals or agents (but also organizations and institutions)
change their adaptive traits over time as a consequence of their
experience. If so, how?
\end{question}
\textit{[Answer here.]}

\subsection{Prediction}
\begin{question}
Prediction is fundamental to successful decision-making; if an agent's
adaptive traits or learning procedures are based on estimating future
consequences of decisions, how do agents predict the future conditions
(either environmental or internal) they will experience? If appropriate,
what internal models are agents assumed to use to estimate future
conditions or consequences of their decisions? What tacit or hidden
predictions are implied in these internal model assumptions?
\end{question}
\textit{[Answer here.]}

\subsection{Sensing}
\begin{question}
What internal and environmental state variables are individuals assumed
to sense and consider in their decisions? What state variables of which
other individuals and entities can an individual perceive; for example,
signals that another individual may intentionally or unintentionally
send? Sensing is often assumed to be local, but can happen through
networks or can even be assumed to be global. If agents sense each other
through social networks, is the structure of the network imposed or
emergent? Are the mechanisms by which agents obtain information modeled
explicitly, or are individuals simply assumed to know these variables?
\end{question}
\textit{[Answer here.]}

\subsection{Interaction}
\begin{question}
What kinds of interactions among agents are assumed? Are there direct
interactions in which individuals encounter and affect others, or are
interactions indirect, e.g., via competition for a mediating resource?
If the interactions involve communication, how are such communications
represented?
\end{question}
\textit{[Answer here.]}

\subsection{Stochasticity}
\begin{question}
What processes are modeled by assuming they are random or partly random?
Is stochasticity used, for example, to reproduce variability in
processes for which it is unimportant to model the actual causes of the
variability? Is it used to cause model events or behaviors to occur with
a specified frequency?
\end{question}
\textit{[Answer here.]}

\subsection{Collectives}
\begin{question}
Do the individuals form or belong to aggregations that affect, and are
affected by, the individuals? Such collectives can be an important
intermediate level of organization in an ABM; examples include social
groups, fish schools and bird flocks, and human networks and
organizations. How are collectives represented? Is a particular
collective an emergent property of the individuals, such as a flock of
birds that assembles as a result of individual behaviors, or is the
collective simply a definition by the modeler, such as the set of
individuals with certain properties, defined as a separate kind of
entity with its own state variables and traits?
\end{question}
\textit{[Answer here.]}

\subsection{Observation}
\begin{question}
What data are collected from the ABM for testing, understanding, and
analyzing it, and how and when are they collected? Are all output data
freely used, or are only certain data sampled and used, to imitate what
can be observed in an empirical study (``Virtual Ecologist'' approach;
Zurell et al.\ 2010)?
\end{question}
\textit{[Answer here.]}

% =====================================================================
% CATEGORY: Details
% =====================================================================

\section{Initialization}
\begin{question}
What is the initial state of the model world, i.e., at time $t = 0$ of
a simulation run? In detail, how many entities of what type are there
initially, and what are the exact values of their state variables (or
how were they set stochastically)? Is initialization always the same, or
is it allowed to vary among simulations? Are the initial values chosen
arbitrarily or based on data? References to those data should be
provided.
\end{question}
\textit{[Answer here.]}% Official explanation: model results cannot be accurately replicated
% unless the initial conditions are known; different analyses can depend
% quite differently on initial conditions.

% =====================================================================

\section{Input data}
\begin{question}
Does the model use input from external sources such as data files or
other models to represent processes that change over time?
\end{question}
\textit{[If the model does NOT use input data, you MUST nevertheless
include this element with the exact statement: ``The model does not use
input data to represent time-varying processes.'' Otherwise describe
each input: what it represents, its units, its source, how it was or can
be obtained, any preparation methods, and uncertainties/biases. If input
was generated by another model, document how.]}
% Official note: 'Input data' does NOT refer to parameter values or
% initial values of state variables (those go to Elements 5 and 7).

% =====================================================================

\section{Submodels}
\begin{question}
What, in detail, are the submodels that represent the processes listed
in `Process overview and scheduling'? What are the model parameters,
their dimensions, and reference values? How were submodels designed or
chosen, and how were they parameterized and then tested?
\end{question}
% Official rules: present the factual description first — equation(s)
% and algorithms — clearly separated from additional information
% (origin of the submodel, why it was chosen, parameterization).
% Parameter definitions, units and values should be presented in
% TABLES. Justification should include references to relevant
% literature and evidence of independent implementation, testing,
% calibration and analysis of submodels.

\subsection{Submodel 1: \textit{[name matching a process of Element 3]}}
\textit{[Equations / algorithms first, then origin and rationale,
parameter table:]}

\begin{center}
\begin{tabular}{lllll}
\textbf{Parameter} & \textbf{Meaning} & \textbf{Unit} & \textbf{Value} & \textbf{Source} \\
\hline
 & & & & \\
 & & & &
\end{tabular}
\end{center}

\subsection{Submodel 2: \textit{[name]}}
\textit{[As above. Add one subsection per submodel listed in Element 3.]}

% Official closing note: ODD-based model descriptions consist of the
% seven elements above; in most cases a "Simulation experiments" (or
% model analysis) section should FOLLOW the model description.

% \section{Simulation experiments} % uncomment if needed
% \textit{[Describe the simulation experiments/analyses performed with
% the model: scenarios, treatments, replication, outputs measured.]}

\begin{thebibliography}{9}
\bibitem{Grimm2006}
Grimm, V., Berger, U., Bastiansen, F., et al. (2006)
A standard protocol for describing individual-based and agent-based models.
\emph{Ecological Modelling}, 198, 115--126.
DOI: 10.1016/j.ecolmodel.2006.04.023

\bibitem{Grimm2010}
Grimm, V., Berger, U., DeAngelis, D.L., Polhill, J.G., Giske, J.,
Railsback, S.F. (2010)
The ODD protocol: A review and first update.
\emph{Ecological Modelling}, 221(23), 2760--2768.
DOI: 10.1016/j.ecolmodel.2010.08.019

\bibitem{Railsback2001}
Railsback, S.F. (2001)
Conceptual design and simulation analysis of individual-based models.
In: \emph{Proceedings of the Third International Symposium on Tools and
Methods of Competitive Engineering}, Delft University Press, pp.~1099--1112.

\bibitem{GrimmRailsback2005}
Grimm, V. and Railsback, S.F. (2005)
\emph{Individual-Based Modeling and Ecology}. Princeton University Press,
Princeton, NJ.

\bibitem{Zurell2010}
Zurell, D., Jeltsch, F. and Grimm, V. (2010)
Ecological modelling and data availability: do we really need all of
these? \emph{Journal of Applied Ecology}, 47, 623--626.
\end{thebibliography}

\end{document}
